\documentclass[12pt,a4paper,openany]{report}

\usepackage{indentfirst}
\usepackage{xurl}

\usepackage{amsmath}
\usepackage{titlesec}

% Заголовки 14 pt, основной текст 12 pt
\newcommand{\HeadingFont}{\normalfont\fontsize{14}{17}\selectfont\bfseries}

\titleformat{\chapter}[hang]
  {\HeadingFont\centering}
  {\thechapter}{1em}{}
\titleformat{name=\chapter,numberless}[hang]
  {\HeadingFont\centering}
  {}{0pt}{}
\titlespacing*{\chapter}{0pt}{0pt}{12pt}

\titleformat{\section}
  {\HeadingFont}
  {\thesection}{1em}{}
\titleformat{name=\section,numberless}
  {\HeadingFont}
  {}{0pt}{}
\titlespacing*{\section}{0pt}{12pt}{6pt}

\usepackage[T2A]{fontenc}
\usepackage[russian]{babel}

\usepackage[
    left=3cm,
    right=1cm,
    top=2cm,
    bottom=2cm
]{geometry}

\usepackage[hidelinks]{hyperref}

\setlength{\parindent}{1.25cm}
\linespread{1.3}

% Не разрывать абзац между страницами
\interlinepenalty=10000
\clubpenalty=10000
\widowpenalty=10000
\displaywidowpenalty=10000
\brokenpenalty=10000

% Не вылезать за правое поле (1 см)
\emergencystretch=2.5em
\tolerance=2500
\hyphenpenalty=50
\exhyphenpenalty=50
\Urlmuskip=0mu plus 3mu
\makeatletter
\g@addto@macro\UrlBreaks{\do\/\do\-\do\.\do\_\do\~\do\#}
\makeatother

\begin{document}
\pagenumbering{gobble}

\begin{titlepage}
\centering

\textbf{МИНИСТЕРСТВО ПРОСВЕЩЕНИЯ РОССИЙСКОЙ ФЕДЕРАЦИИ\\[0.5cm]
Федеральное государственное бюджетное образовательное учреждение высшего образования\\
«Херсонский государственный педагогический университет»}\\[0.5cm]
Институт информационных технологий\\[0.5cm]
Кафедра программной инженерии и информационных технологий

\vfill

\textbf{КУРСОВАЯ РАБОТА}\\
\textbf{на тему: «Разработка сайта-визитки\\
фуллстек-разработчика»}\\[0.5cm]
\textbf{по дисциплине «Современные операционные системы\\
и базы данных»}

\vfill

\begin{flushleft}
Обучающегося 2 курса группы 12-25РПм Направления подготовки 09.04.04\\
Трифонова Д.С.\\[0.5cm]

Научный руководитель:\\
Галлини Н.И.
\end{flushleft}

\vfill

Херсонская область - 2026

\end{titlepage}

\clearpage
\pagenumbering{arabic}

\tableofcontents
\clearpage

% начало текста для антиплагиата
%<*plagiarism>

\chapter*{ВВЕДЕНИЕ}
\addcontentsline{toc}{chapter}{ВВЕДЕНИЕ}

В профессиональной среде первого впечатления часто достаточно нескольких секунд просмотра персональной страницы. Прежде чем состоится разговор или переписка, заинтересованная сторона уже успевает оценить, как специалист оформляет своё присутствие в сети, насколько ясно изложены сильные стороны и насколько удобно начать контакт. Компактная персональная страница, которую принято называть визиткой, как раз и решает эту задачу: она собирает ключевые сведения в одном экране и задаёт визуальный тон.

Если речь идёт о фуллстек-разработчике, ценность такого ресурса возрастает. Страница одновременно показывает практический навык вёрстки и клиентской логики, аккуратность в подготовке к выкладке и умение держать интерфейс понятным на разных устройствах. Даже когда объём контента невелик, приходится продумать иерархию блоков, скорость открытия, устойчивость вёрстки и порядок файлов на сервере.

Тема курсовой работы актуальна потому, что веб-страница существует не «сама по себе», а внутри связки «браузер - сеть - операционная среда хостинга». Для дисциплины «Современные операционные системы и базы данных» здесь важны как минимум три аспекта: как ОС обслуживает выдачу файлов, где физически лежат данные и чем клиентское хранилище отличается от серверной СУБД. Отдельно учитывается подготовка ресурса к индексации.

Цель работы - создать и подробно описать персональную визитку фуллстек-разработчика, пояснив выбранный стек и схему размещения в сети.

Чтобы достичь цели, последовательно решаются следующие задачи:
\begin{itemize}
\item выяснить, какие ожидания предъявляются к персональной веб-визитке специалиста;
\item разобрать, как браузер и сервер взаимодействуют при открытии страницы и какую роль здесь играет ОС;
\item сопоставить варианты хранения сведений: файлы на диске, браузерное хранилище и классическая СУБД;
\item спроектировать и собрать одностраничный интерфейс с адаптацией и умеренной интерактивностью;
\item выложить результат на хостинг и оформить минимальный набор средств поисковой подготовки.
\end{itemize}

Объект исследования - создание небольшого персонального веб-ресурса специалиста.

Предмет исследования - приёмы построения сайта-визитки фуллстек-разработчика с опорой на свойства операционных систем и способы хранения данных.

В работе применяются изучение источников и документации, сопоставление вариантов архитектуры, макетирование интерфейса, практическая реализация, проверка в браузерах и публикация на хостинге.

Практический итог - действующая визитка по адресу \url{http://trifonix.beget.tech/}. Полученные решения можно переносить на другие лёгкие персональные страницы.


\chapter{ТЕОРЕТИЧЕСКИЕ ОСНОВЫ РАЗРАБОТКИ САЙТА-ВИЗИТКИ}

\section{Сайт-визитка как форма цифрового представления специалиста}

Под сайтом-визиткой понимают небольшой публичный ресурс, где человек или организация кратко представляются аудитории~[1]. У него нет задач полноценного портала или торговой площадки: не нужны корзина, личный кабинет и сложные сценарии. На первый план выходят ясная подача, быстрый отклик страницы и понятный путь к контакту.

В персональной визитке разработчика обычно встречаются:
\begin{itemize}
\item блок идентичности: ФИО, снимок, должность или роль;
\item короткое резюме умений;
\item сведения о местах работы и учёбе;
\item список технологий;
\item выходы на проекты и мессенджеры.
\end{itemize}

С инженерной позиции такой ресурс ближе к информационной витрине, чем к транзакционной системе. Контент меняется редко, поэтому его удобно держать в разметке. Вместе с тем на клиенте допустимы небольшие сценарии - смена оформления, подсчёт стажа, анимация акцентов. Главное остаётся прежним: показать сведения, а не обслуживать поток операций.

Для фуллстек-разработчика визитка ещё и «пробный стенд» полного пути поставки: сверстать, оформить стилями, добавить логику, положить файлы на сервер, проверить служебные \path{robots.txt} и \path{sitemap.xml}, убедиться в корректном виде на телефоне и мониторе. В итоге страница работает и как продукт, и как фрагмент портфолио.

\section{Клиент-серверная модель и роль операционной системы}

Сеть веб-ресурсов опирается на разделение ролей клиента и сервера~[2]. Клиент - это браузер: он отправляет запрос по HTTP или HTTPS, получает ответ, разбирает разметку, применяет таблицы стилей, запускает сценарии и рисует интерфейс. Сервер принимает обращение, находит нужный файл либо формирует ответ программно и возвращает его обратно.

ОС на стороне сервера создаёт условия для работы веб-сервиса (часто Apache либо Nginx): поднимает процессы, контролирует права на каталоги, открывает сетевые порты, ведёт журналы, обслуживает сертификаты~[3]. Даже если страница полностью статическая, без корректной настройки ОС нельзя гарантировать безопасный доступ к публичным файлам и стабильную выдачу.

На устройстве пользователя своя ОС (Windows, macOS, Android, iOS и другие) запускает браузер как приложение, выделяет ему память и сеть, ограничивает доступ к системе. Внутри браузера действуют собственные правила: исполнение JavaScript преимущественно в одном потоке интерфейса, политика same-origin, лимиты на объём \texttt{localStorage}. Эти ограничения нужно учитывать уже на этапе проектирования.

Открытие визитки обычно проходит так:
\begin{enumerate}
\item человек набирает адрес или переходит по ссылке;
\item служба имён преобразует домен в IP;
\item устанавливается защищённое соединение;
\item сервер под управлением ОС отдаёт \texttt{index.html} и связанные файлы;
\item браузер строит дерево документа, накладывает стили и выполняет скрипты.
\end{enumerate}

Зная эту цепочку, проще выбрать площадку размещения: shared-хостинг, VPS или чисто статический хост. В курсовом проекте использован виртуальный хостинг Beget: ОС и веб-сервер предоставляет площадка, а автор отвечает за наполнение открытого каталога.

\section{Хранение данных в веб-приложениях}

СУБД нужны там, где сведения нужно систематически сохранять, искать и изменять~[4]. В типичном динамическом сайте серверная база (PostgreSQL, MySQL и аналоги) держит пользователей, материалы и события; приложение читает и пишет её через программный слой, а ОС обеспечивает процессы СУБД и доступ к дискам.

Для визитки отдельный серверный склад часто избыточен: тексты и изображения зафиксированы и могут жить прямо в файлах проекта. При этом хранение всё равно проявляется в двух плоскостях.

На сервере файловая система ОС хранит страницу, иллюстрации, иконки, манифест и SEO-служебники. Здесь важны порядок каталогов, типы содержимого и правила кэша.

В браузере доступен Web Storage. Через \texttt{localStorage} можно оставлять строковые пары ключ-значение между визитами~[5]. В отличие от cookie, эти записи не уезжают на сервер с каждым запросом. Реляционных выборок, транзакций и совместного доступа нескольких пользователей там нет, зато удобно помнить личные настройки оформления.

В выполненном проекте в \texttt{localStorage} фиксируется выбранная палитра (светлая либо тёмная). Так разделяются:
\begin{itemize}
\item содержимое самой визитки, которое сопровождает исходный код;
\item предпочтения посетителя, привязанные к его браузеру.
\end{itemize}

Отдельно стоит отметить машиночитаемые описания Open Graph и JSON-LD. Это не замена СУБД, но способ сообщить поисковикам и мессенджерам, о ком страница и какое изображение к ней относится.


\chapter{ПРОЕКТИРОВАНИЕ САЙТА-ВИЗИТКИ}

\section{Требования к функциональности и качеству}

Перед сборкой интерфейса зафиксирован перечень ожидаемых возможностей:
\begin{enumerate}
\item показать фото, ФИО, роль и короткое описание умений;
\item вывести места работы и образование;
\item перечислить ключевые технологии;
\item дать анонс собственного SaaS и ссылку на него;
\item обеспечить быстрый выход на почту, GitHub и Telegram;
\item позволить менять светлую и тёмную палитру с запоминанием выбора;
\item автоматически считать срок работы «по настоящее время»;
\item уложить материал в макет «экран смартфона» так, чтобы основной каркас не требовал прокрутки.
\end{enumerate}

К качеству предъявлялись требования лаконичного внешнего вида, быстрого открытия, понятных элементов управления, готовности к выкладке и базовой поисковой подготовки (title, description, canonical, Open~Graph, sitemap).

Условие «один экран» заставило собрать сведения мозаикой внутри рамки телефона. Такой каркас усиливает ощущение мобильного формата и делает композицию собранной.

\section{Выбор технологического стека}

Реализация опирается на HTML5, CSS3 и JavaScript без серверных каркасов. Выбор следует из характера задачи: тексты почти не меняются, реакции происходят локально, а выкладка на обычный хостинг должна оставаться простой.

Разметка отвечает за смысл страницы: главный заголовок с именем, заголовки блоков, контактные ссылки и служебные метаданные. Стили строят сетку, неоновую эстетику, темы через CSS-переменные и движения элементов. Сценарии отвечают за стартовый экран загрузки, подсчёт стажа, смену темы и лёгкий параллакс фона.

Отказ от серверной СУБД на текущем шаге не уводит работу в сторону от дисциплины. Напротив, в тексте явно показано:
\begin{itemize}
\item где хватает файлов под контролем ОС;
\item где уместно браузерное хранилище;
\item когда без полноценной базы уже не обойтись (учётки, CMS, комментарии, поток событий).
\end{itemize}

В перспективе можно добавить API и таблицу проектов, чтобы обновлять сведения без правки HTML. Для учебной цели хватает статической поставки, и она проще для демонстрации идей курса.

\section{Информационная архитектура и интерфейс}

Страница собрана из одной мозаики блоков:
\begin{itemize}
\item верхняя полоса - личность, статус, контакты;
\item опыт - две актуальные позиции с датами;
\item образование - вуз, направление, ссылка на сайт;
\item навыки - набор компактных «плашек»;
\item SaaS - краткий анонс DigiGu.
\end{itemize}

Визуально всё помещено в корпус смартфона с узнаваемыми деталями: зона Dynamic Island и нижняя полоска Home. В «островке» расположена кнопка «ТЕМА», поэтому смена оформления находится на видном месте.

Цвет помогает читать карту блоков: опыт подсвечен розовым, учёба - синим, навыки - оранжевым, SaaS - зелёным. Даты двух мест работы дополнительно различаются песочным и виноградным тонами, а сами карточки слегка тонированы фоном.

На узких экранах media-запросы меняют сетку навыков, укорачивают длинные подписи и перестраивают шапку. Так выдерживается логика mobile-first.


\chapter{ПРАКТИЧЕСКАЯ РЕАЛИЗАЦИЯ И РАЗВЁРТЫВАНИЕ}

\section{Структура проекта и основные компоненты}

Собранный результат - набор статических файлов:
\begin{itemize}
\item \texttt{index.html} содержит разметку, стили и клиентские сценарии;
\item \texttt{scr/photo/me.jpg} хранит портрет;
\item каталог \texttt{scr/fav/} включает фавиконки и web-manifest;
\item \texttt{robots.txt}, \texttt{sitemap.xml}, \texttt{.htaccess}, \texttt{404.html} обслуживают публикацию.
\end{itemize}

Соединение представления и логики в одном \texttt{index.html} ускоряет перенос: достаточно скопировать каталог в публичную область хостинга. Для учебного объёма это приемлемо; в промышленной практике стили и скрипты чаще выносят и собирают отдельно.

Среди заметных решений:
\begin{enumerate}
\item палитры через CSS-переменные и атрибут \texttt{data-theme};
\item запоминание темы в \texttt{localStorage} ключом \texttt{kyr2-theme};
\item подсчёт стажа из \texttt{data-start} формата \texttt{MM.YYYY} с русскими формами слов;
\item поочерёдная подсветка навыков, которая замирает при наведении;
\item SEO-описания и JSON-LD сущностей \texttt{Person} и \texttt{WebPage}.
\end{enumerate}

\section{Взаимодействие с операционной средой хостинга}

Ресурс доступен по адресу \url{http://trifonix.beget.tech/}. Площадка Beget даёт окружение Linux и Apache. Через \path{.htaccess} заданы:
\begin{itemize}
\item стартовый документ \texttt{index.html};
\item базовые security-заголовки;
\item сжатие текстовых ответов;
\item сроки кэша для картинок и HTML.
\end{itemize}

С позиции ОС это настройка уровня пользователя виртуального хоста: ядро системы не администрируется, но правила раздачи файлов контролируются. Подобная модель обычна для shared- и облачных площадок.

Файл \path{robots.txt} сообщает роботам границы обхода, а \path{sitemap.xml} указывает канонический адрес главной. Вместе с мета-описаниями это укрепляет техническую готовность страницы к индексации.

\section{Клиентское хранение настроек и вычисляемые данные}

Смена темы не ходит на сервер. При открытии страницы сценарий читает запись из \texttt{localStorage} и применяет палитру до показа основного интерфейса, чтобы уменьшить вспышку «чужого» цвета. Нажатие «ТЕМА» обновляет и экран, и сохранённое значение.

Срок работы считается в браузере по системной дате устройства. Для верхней карточки взята точка старта 04.2026, для нижней - 03.2025. Подпись выводится под «н.в.» двумя строками: годы и месяцы. Так демонстрируется работа со временем без СУБД: исходные даты лежат в разметке, а результат появляется в момент открытия.

Методически это подчёркивает простую мысль: не каждое вычисляемое поле обязано жить в таблице. Атрибуты HTML и клиентский расчёт уместны, пока нет нужды синхронно менять сведения для множества посетителей из одного центра.

\section{Тестирование и результаты}

Контроль включал:
\begin{itemize}
\item просмотр на широком и узком экране;
\item проверку смены темы и её сохранения после обновления;
\item сверку подсчёта стажа;
\item работоспособность контактов и кнопки темы;
\item наличие иконок вкладки и открытие боевого URL.
\end{itemize}

В итоге получена рабочая визитка с актуальными сведениями об авторе. Она фиксирует практику вёрстки, клиентских сценариев и выкладки и может служить и портфолио-элементом, и заготовкой для следующих персональных страниц.

%*</plagiarism>

\chapter*{ЗАКЛЮЧЕНИЕ}
\addcontentsline{toc}{chapter}{ЗАКЛЮЧЕНИЕ}

В курсовой работе рассмотрена разработка сайта-визитки фуллстек-разработчика в контексте современных операционных систем и подходов к хранению данных. Основные результаты состоят в следующем.

\begin{enumerate}
\item Показано, что сайт-визитка является эффективным средством профессиональной самопрезентации и может служить демонстрацией компетенций веб-разработчика.

\item Установлено, что даже статический веб-ресурс опирается на клиент-серверную модель и сервисы операционной системы хостинга: файловую систему, веб-сервер, сетевой стек и политики доступа.

\item Проанализированы варианты хранения данных. Для контента визитки достаточно файлового размещения; для пользовательских настроек интерфейса целесообразно клиентское хранилище \texttt{localStorage}; серверная СУБД требуется при появлении многопользовательской динамики и централизованного управления контентом.

\item Спроектирован и реализован одностраничный адаптивный сайт с блоками опыта, образования, навыков и проекта, а также с переключением темы и автоматическим расчётом стажа.

\item Выполнено развёртывание на хостинге Beget по адресу \url{http://trifonix.beget.tech/}, подготовлены файлы SEO и конфигурации веб-сервера.
\end{enumerate}

Практическая значимость подтверждается наличием опубликованного результата. Дальнейшее развитие может включать серверный API, базу данных проектов, многоязычность и сбор обезличенной статистики посещений.


\chapter*{Список использованных источников}
\addcontentsline{toc}{chapter}{Список использованных источников}

\begin{enumerate}
\item Никсон Р. Создаём динамические веб-сайты с помощью PHP, MySQL, JavaScript, CSS и HTML5. - СПб.: Питер, 2021. - 816 с.
\item Таненбаум Э., Бос Х. Современные операционные системы. - 4-е изд. - СПб.: Питер, 2020. - 1120 с.
\item Олифер В. Г., Олифер Н. А. Компьютерные сети. Принципы, технологии, протоколы. - СПб.: Питер, 2020. - 1008 с.
\item Дейт К. Дж. Введение в системы баз данных. - М.: Вильямс, 2019. - 1328 с.
\item Документация MDN Web Docs. Window.localStorage [Электронный ресурс]. - Режим доступа: \url{https://developer.mozilla.org/ru/docs/Web/API/Window/localStorage} (дата обращения: 10.09.2026).
\item Документация MDN Web Docs. HTML: HyperText Markup Language [Электронный ресурс]. - Режим доступа: \url{https://developer.mozilla.org/ru/docs/Web/HTML} (дата обращения: 10.09.2026).
\item Спецификация Schema.org. Person [Электронный ресурс]. - Режим доступа: \url{https://schema.org/Person} (дата обращения: 10.09.2026).
\item Хостинг Beget. Документация и справка [Электронный ресурс]. - Режим доступа: \url{https://beget.com/ru/kb} (дата обращения: 10.09.2026).
\item Гамма Э., Хелм Р., Джонсон Р., Влиссидес Дж. Приёмы объектно-ориентированного проектирования. Паттерны проектирования. - СПб.: Питер, 2020. - 366 с.
\item Фримен Э., Робсон Э. Изучаем HTML, CSS и JavaScript. - СПб.: Питер, 2021. - 720 с.
\end{enumerate}

\end{document}
